% Figure: pointwise trap -- a piecewise ±x counterexample.
% Source for Figure~\ref{fig:trap}. TikZ only; regenerable by compiling main.tex.
\begin{tikzpicture}[
  lab/.style={font=\footnotesize\sffamily,text=slate,align=center},
  dot/.style={circle,fill=accent,inner sep=1.6pt},
]
\draw[rulegray,->] (-2.4,0) -- (2.4,0);
\draw[rulegray,->] (0,-2.2) -- (0,2.2);
% f(x)=x on x<=0 part drawn solid right-half, f(x)=-x on x>0:
\draw[ink,thick] (-2,2) -- (0,0);
\draw[ink,thick] (0,0) -- (2,-2);
\draw[accent,dashed,thick] (-2,-2) -- (0,0);
\draw[accent,dashed,thick] (0,0) -- (2,2);
\node[lab] at (2.15,1.9) {$+x$};
\node[lab] at (2.15,-1.9) {$-x$};
\node[lab,align=center]{each point obeys $f(t)=\pm t$ --- the sign\\
may change by point, so the global claim fails};
\end{tikzpicture}
